\documentclass[10pt]{article}
\usepackage[margin=20mm]{geometry}
\usepackage{amsmath,amssymb,mathtools,fontspec,unicode-math,xurl,hyperref}
\IfFontExistsTF{DejaVu Serif}{\setmainfont{DejaVu Serif}}{\setmainfont{Latin Modern Roman}}
\setmathfont{latinmodern-math.otf}
\setlength{\parindent}{0pt}\setlength{\parskip}{6pt}\sloppy
\title{Initial scientific review of
Deformation theory and partition Lie algebras}
\author{}
\date{}
\begin{document}
\maketitle
Abstract: The inspected arguments preserve the connectivity, completeness and finite-type restrictions needed for the formal-moduli classification. One Major issue is established in the printed derivation of the many-generator formula: its binomial expansion removes a nontrivial symmetric-group action, and a subsequent identity for general inputs loses tensor factors. Explicit finite-dimensional counterexamples are given. This is a repairable proof issue, not a counterexample to the stated classification or basis theorems. Coverage and remaining dependencies are specified below.\par
\section*{Identity, version and scope}
This report concerns M25-0047, D. Lukas B. Brantner and Akhil Mathew, Deformation theory and partition Lie algebras, Acta Mathematica 235 (2025), 1–148, DOI 10.4310/acta.2025.v235.n1.a1 [BM]. The assigned 148-page publisher PDF was reviewed; printed and PDF page numbers coincide. Its SHA-256 is\par
`f01a1a4ab4d8e5fbb2dff271f68beb47`
`ec835278a696b5f150d021762648254c`.\par
The later arXiv v5, dated 27 November 2025, was consulted for a limited comparison; it does not replace the publisher file.\par
The review exposed 82 of 148 pages as text: 1–18, 24–29, 44–50, 59–68, 77–91, 96–101, 103–114, 121–123 and 130–134. Twenty-four source pages were also inspected as rendered images. Full extraction and theorem indexing served navigation; neither constitutes review of unread pages. Targets were the classification hypotheses, comonadicity and completion, the relative-base reduction, explicit partition-complex formulas, and selected basis-calculation steps. This is a systematic claim-focused initial review, not a line-by-line verification of all 148 pages.\par
\section*{Problem M01: the many-generator derivation loses equivariant information}
Severity: Major, confined to a demonstrated core proof calculation.\par
\subsection*{Exact locations and dependency}
On p. 132, the proof of Proposition 7.40 first gives a binomial expansion for finite pointed \(X\), finite \(V\), and arbitrary integers \(\ell_i\), then extends it to (7.7). It pulls suspension factors outside subgroup coinvariants while leaving the permutation action on \(V^{\otimes n}\) untwisted. The last display on p. 132, repeated after dualisation on p. 133, makes a further assertion for arbitrary \(V\). The proof of Theorems 7.3 and 7.4 ends on p. 134 by using Proposition 7.40 for every number of generators greater than one. These are the advertised basis theorems, restated as Theorems 1.18 and 1.22. The invalid calculation thus has an explicit dependency on a core contribution.\par
\subsection*{Complete counterexample to the binomial expansion}
Equation reference: signed. Take \[p=3,\quad n=2,\quad m=1,\quad
 \ell_1=1,\quad V=\mathbb F_3,\quad X=S^0\] with trivial \(\Sigma_2\)-action on the single non-basepoint. All the displayed finiteness and integrality assumptions hold. Put \(W=\Sigma^{-1}\mathbb F_3\). The transposition on \(W^{\otimes2}\) acts by \(-1\), since both factors have odd degree. Hence \[
 \mathbb F_3[X]\otimes_{\Sigma_2} W^{\otimes2}=0.\] The coinvariant relation identifies a generator \(z\) with \(-z\), so \(2z=0\), and 2 is invertible. There are no hidden derived higher coinvariants: averaging makes invariants and coinvariants exact when the group order is invertible. Homotopy orbits give the same answer.\par
Equation reference: untwisted. The printed right-hand side has only \(n_1=2\), and gives \[
 \Sigma^{-2}\bigl(\mathbb F_3[X]\otimes_{\Sigma_2}V^{\otimes2}\bigr)
 =\Sigma^{-2}\mathbb F_3\ne0.\] Thus the preliminary expansion is false. Formula (7.7) has the same problem: in this semisimple arity, \(F_{S^0}\) is the second symmetric-power summand on complexes, so \[F_{S^0}(\Sigma^{-1}\mathbb F_3)=0,\qquad
 \Sigma^{-2}F_{S^0}(\mathbb F_3)=\Sigma^{-2}\mathbb F_3.\] This is not a distinction between strict and homotopy coinvariants.\par
A valid binomial expansion retains \[\bigotimes_i(\Sigma^{-\ell_i}V)^{\otimes n_i}\] inside the coinvariants for \(H=\prod_i\Sigma_{n_i}\). For discrete \(V\), pulling the total suspension outside must retain the character \(\bigotimes_i\operatorname{sgn}_{n_i}^{\ell_i}\) on \(V^{\otimes n}\). Derived extensions must carry the corresponding equivariant suspension data; an untwisted scalar shift cannot replace it.\par
\subsection*{A second check excludes a correction of only one sign}
Write \(C_r=\Sigma|\Pi_r|^\diamond\), with the intended convention \(C_1=S^0\). The last display on p. 132 uses the exponent \(1+\sum_i(1-\ell_i)n_i/d\), whereas the desired dual input requires \(-1+\sum_i(1-\ell_i)n_i/d\). Changing that sign alone does not make the assertion true for arbitrary \(V\).\par
Work over \(\mathbb F_5\), take \(n=2\), \(m=2\), \(\ell_1=\ell_2=0\), and \(V=\mathbb F_5^2\). Inspect mixed multidegree \((n_1,n_2)=(1,1)\). Since \(C_2=S^1\) with trivial \(\Sigma_2\)-action, the left-hand mixed summand is \[\Sigma(V\otimes V),\qquad \dim_{\mathbb F_5}=4.\] This is also the mixed term of \(\Sigma\operatorname{Sym}^2(V\oplus V)\). The right has only \(d=1\) and the Lyndon word \(x_1x_2\). Even after correcting the suspension sign, its term is \[F_{C_1}(\Sigma V)=\Sigma V,\qquad \dim_{\mathbb F_5}=2.\] A tensor power has disappeared. This example is in nonnegative degrees and semisimple characteristic, independently of the odd-shift example. The general-input formula must retain the appropriate \(V^{\otimes n/d}\), or specialise to \(V=\mathbb F_p\) before the simplification.\par
\subsection*{Counterevidence, repair and restricted impact}
There is strong counterevidence to interpreting this as a false main theorem. Proposition 7.40 is stated for sums of shifts of the one-dimensional \(\mathbb F_p\); it does not assert the false arbitrary-\(V\) identity. Its exponent \(1+\sum_i(\ell_i-1)|w|_i\) has the expected bracket degree. Theorem 7.39 retains permutation spheres and provides the correct starting point.\par
After specialising \(V=\mathbb F_p\), retain all shifted factors in the \(H\)-equivariant binomial expansion. For \(d\mid\gcd(n_i)\), restriction to the diagonal \(\Sigma_d\) groups these as the \(d\)-fold tensor power of \[\bigotimes_i(\Sigma^{-\ell_i}\mathbb F_p)^{\otimes n_i/d}.\] The sphere \((S^{n/d-1})^{\wedge d}\) in Theorem 7.39 combines with them to give the \(d\)-fold tensor power of a line in degree \[a=n/d-1-\sum_i\ell_i n_i/d
   =-1+\sum_i(1-\ell_i)n_i/d.\] Duality gives \(-a=1+\sum_i(\ell_i-1)n_i/d\), exactly the exponent in the statement of Proposition 7.40. Keeping the grouped tensor factors retains their Koszul action. In the spectral case this is compatible with the bounded-above formula of Proposition 5.35. The strict case requires the corresponding genuine-equivariant and right-left-extension argument, rather than (7.7). A full reconstruction of that extension argument for every shift is not claimed here.\par
The established impact is an invalid step in the multigenerator basis proof, with a clear correction route. Neither a counterexample to Theorems 7.3–7.4 nor damage to the separate classification Theorems 1.12, 1.21 and 6.27 has been established. Major follows the specified standard for a demonstrated invalid derivation used by a core theorem; repairability does not remove that dependency. The later v5 HTML still contains the untwisted expansion and general-\(V\) formula. Limited primary-source checks found no explicit erratum for this issue; novelty is not certified.\par
\section*{Other reasoning and boundary checks}
\subsection*{Formal-moduli hypotheses and completion}
Definitions 1.3 and 1.5 distinguish finite total homotopy for Artinian inputs from degreewise finite type for complete Noetherian objects. The designated maps in the pullback condition must be surjective on \(\pi_0\). Claims about arbitrary nonconnective algebras or arbitrary pullbacks were not substituted.\par
In Proposition 2.11, pp. 17–18, a sufficiently high finite skeleton of a simplicial connective object computes a fixed homotopy degree uniformly across the filtration. Finite colimits preserve complete filtered objects; the Milnor \(\lim\) and \(\lim^1\) terms can then be controlled degree by degree. This does not say arbitrary colimits preserve completeness.\par
In Theorem 4.20, pp. 61–65, a cotangent equivalence gives an associated-graded equivalence, and completeness gives conservativity. Totalisation uses a cotangent-split cosimplicial object and its adic lift, not an unrestricted interchange of cotangent fibres and limits. For square-zero inputs, duality exchanges the relevant pullback with a pushout and the free functor preserves the pushout. The general step uses cobar resolutions, \(\pi_0\)-surjectivity and coherence. The underlying filtered-totalisation lemmas were traced but not all independently reproved.\par
Lemma 5.21 and Proposition 5.22, pp. 77–78, use bounded-below almost-perfect modules and Nakayama’s lemma; the least nonzero homotopy degree is legitimate. The cotangent-space map makes the map on complete local rings surjective, after which the relative cotangent module is almost perfect. The formal-etale and connectivity criteria remain external dependencies. Proposition 5.28, p. 80, uses finite positive filtration degrees and completion of a finitely generated graded module, not an arbitrary direct-sum/product interchange.\par
\subsection*{Relative bases and duality}
Propositions 6.5–6.6 and Theorem 6.7, pp. 97–99, retain derived completeness and connectivity. Perfect approximations to an almost-perfect \(N\) have increasingly connective cofibres, giving degreewise control of tensor product versus totalisation. The Cech nerve splits after tensoring with \(k\), and conservativity on complete connective modules gives descent. This would not justify descent for all unbounded modules.\par
Theorems 6.15 and 6.24, pp. 103–107, reduce to the case over \(k\). Their duality is on degreewise finite coconnective objects. Double duality was not extended to arbitrary infinite-dimensional modules. The formal-moduli identification in Theorem 6.27 also invokes Lurie’s axiomatic deformation theorem; that external theorem was not independently reconstructed.\par
\subsection*{Partition complexes and low weights}
Proposition 5.35 uses a right-left extension for arbitrary \(V\); the simpler homotopy-fixed-point formula requires \(V\) bounded above. Proposition 5.49 instead uses strict fixed points in a cosimplicial model for simplicial commutative rings. These are not interchangeable in positive characteristic.\par
At arity one, the bar model on p. 84 has one non-basepoint simplex with all partitions equal, giving \(T(1)=S^0\). Literal use of the empty proper-partition nerve would be misleading. The primary source [AB, p. 84] explicitly specifies \(C_1=S^0\); the apparent unit discrepancy is therefore resolved by an actual convention and is not a finding.\par
Reviewer-written finite-field chain calculations enumerated the proper-partition posets and all chains for \(2\le n\le5\). Over each of \(\mathbb F_2,\mathbb F_3,\mathbb F_5\), reduced homology is concentrated in degree \(n-3\), with dimensions \(1,2,6,24\). Double suspension gives degrees \(n-1\), as expected. This finite check is not a proof for all arities.\par
At \(p=2\), spectral weight two on a class of degree \(\ell\) is \(\Sigma^{2\ell-1}\mathbb F_2^{hC_2}\), with one class in every degree at most \(2\ell-1\). The basis condition \(i_1\le\ell\), with degree \(\ell+i_1-1\), agrees. For odd primes, the extra self-bracket of an even-degree class has weight two and degree \(2\ell-1\), as required by the shifted Lie convention.\par
The Euler cancellation on pp. 121–122 is \[\sum_{s=r}^{k}(-1)^{s-1}\binom{k-r}{s-r}
 =
 \begin{cases}0,&r<k,\\(-1)^{k-1},&r=k.\end{cases}\] The endpoint simplifications on pp. 122–123 and 133–134 were checked with parity and congruences retained. Independent enumeration tested 249,510 bounded tuples for the former simplification and 2,250 even-generator endpoint cases without mismatch; the bounds are recorded in the script. This does not verify the full Bredon spectral sequence or matching theorem.\par
Corollary 7.9, p. 114, prints even where Proposition 7.8 and the surrounding odd-generator proof require odd degree at odd primes. The context and separate even-generator argument identify the intended correction. This isolated parity-word error and the isolated suspension-sign transcription on pp. 132–133 are not separate substantive findings. M01 rests on the independent action and dimension counterexamples.\par
\section*{Evidence and limitations}
The coverage ledger records the inspected pages and claims. Evidence includes source-page images, a reviewer-written exact-arithmetic script and actual output, source bindings, and compilation records. No author program was supplied in the acquired manifest or executed; the package contains a publisher PDF and metadata rather than a computational repository.\par
Unverified material includes the complete extension calculus in Section 3; all filtered convergence lemmas; the full operadic generalisation in Section 5.3; Appendix A’s hypercoverings and Kan extensions; the full Bredon matching and degeneration argument; the topological-algebra construction of the even-generator cofibration; and the external foundations and every cited theorem. The correction route for M01 is not promoted to a fully verified new proof of the article.\par
The outcome is one supported Major proof-calculation issue with restricted impact, and no additional substantive issue established in the inspected classification arguments. Scientific acceptance remains pending. This report does not certify whole-paper correctness or record author adoption.\par
9 D. L. B. Brantner and A. Mathew, Deformation theory and partition Lie algebras, Acta Math. 235 (2025), 1–148. https://doi.org/10.4310/acta.2025.v235.n1.a1. Assigned publisher version. G. Z. Arone and D. L. B. Brantner, The action of Young subgroups on the partition complex, Publ. Math. IHES 133 (2021), 47–156. https://doi.org/10.1007/s10240-021-00123-7. The arity-one convention on p. 84 was directly consulted. D. L. B. Brantner and A. Mathew, Deformation theory and partition Lie algebras, arXiv:1904.07352v5, 27 November 2025. https://arxiv.org/html/1904.07352v5. Limited comparison of Proposition 7.40; not substituted for the assigned version.\par
\end{document}
